\documentclass[12pt]{article}

\usepackage{preamble}

% ---------------------------------------------------------------------------
% Document-specific macros
% ---------------------------------------------------------------------------
\newcommand{\anote}[1]{\textcolor{red}{\small\textbf{(Ashwin:} #1\textbf{)}}}
% \todo{...} prints the note inline and adds a numbered, hyperlinked entry to
% \listoftodos (written to main.tdo, like the table of contents).
\makeatletter
\newcounter{todo}
\newcommand{\todo}[1]{%
  \refstepcounter{todo}%
  \addcontentsline{tdo}{todo}{\protect\numberline{\thetodo}#1}%
  \textcolor{\mysubcolor}{\small\textbf{[TODO~\thetodo:} #1\textbf{]}}}
\newcommand{\l@todo}{\@dottedtocline{1}{0em}{2em}}
\newcommand{\listoftodos}{%
  \section*{List of TODOs}%
  \@starttoc{tdo}}
\makeatother

\newcommand{\alg}{\mathcal{A}}             % finite-dimensional C*-algebra
\newcommand{\hilbH}{\mathcal{H}}            % Hilbert space
\newcommand{\hilbK}{\mathcal{K}}             % dilation space
\newcommand{\hilbN}{\mathcal{N}}             % null space of a semidefinite form
\newcommand{\bdd}[1]{\mathcal{B}\br{#1}}   % bounded operators
\newcommand{\Mn}[1][n]{\mathbb{M}_{#1}}    % matrix algebra: \Mn = M_n, \Mn[k] = M_k
\newcommand{\algtensor}{\odot}             % algebraic tensor product
\newcommand{\op}{^{\mathrm{op}}}           % opposite algebra
\newcommand{\ip}[2]{\left\langle #1 , #2 \right\rangle}
\newcommand{\ipphi}[2]{\left\langle #1 , #2 \right\rangle_{\varphi}}
\newcommand{\Ltwo}{L^2\br{\alg}}
\newcommand{\vect}[1]{\widehat{#1}}
\newcommand{\Ltn}{\hilbH \tensor \Ltwo}
\newcommand{\Fphi}{F_{\varphi}}
\newcommand{\Choi}{C_{\varphi}}
\newcommand{\comm}{\mathcal{C}}            % commutant rho(A)'
\newcommand{\Esc}{E_{\complex 1}}          % expectation onto the scalars
\newcommand{\ran}{\operatorname{ran}}

% ---------------------------------------------------------------------------
\title{\textbf{Stinespring's Theorem via the Choi Matrix:\\
an \texorpdfstring{\(L^2\)}{L2} Route in Finite Dimensions}%
%\thanks{Acknowledgements common to all authors.}
}

\author{
  Bert Sun \\
  Institute for Quantum Computing \\
  University of Waterloo, Waterloo, Ontario, Canada \\
  \texttt{bert.sun@uwaterloo.ca}
}

\date{\today}

\begin{document}

\maketitle

\begin{abstract}
\todo{abstract}
\end{abstract}

\tableofcontents

% ===========================================================================
\section{Preliminaries}
\label{sec:prelim}

% ---------------------------------------------------------------------------
\subsection{Conventions}
\label{sec:conventions}

% ---------------------------------------------------------------------------
\subsubsection{Sesquilinear forms and bra-ket notation}
\label{sec:conv-braket}

We use two notations for the inner product on a Hilbert space \(\hilbH\), told
apart by the separator.
\begin{itemize}
  \item \emph{Comma (mathematics convention).} \(\ip{x}{y}\) is linear in the
    first argument and conjugate-linear in the second.
  \item \emph{Vertical bar (physics convention).} \(\braket{x}{y}\) is
    conjugate-linear in the first argument and linear in the second.
\end{itemize}
The two are related by
\begin{equation}
\label{eq:bar-comma}
  \braket{x}{y} = \ip{y}{x}, \qquad x, y \in \hilbH .
\end{equation}
A ket \(\ket{x}\) is the vector \(x\), and a bra \(\bra{y}\) is the functional
\(x \mapsto \braket{y}{x} = \ip{x}{y}\). Matrix elements read
\(\bra{y}T\ket{x} = \braket{y}{Tx} = \ip{Tx}{y}\). For \(u, v \in \hilbH\) the
outer product \(\ketbra{u}{v}\) is the rank-one operator
\begin{equation}
\label{eq:ketbra}
  \ketbra{u}{v}\, w \eqdef \braket{v}{w}\, u = \ip{w}{v}\, u ,
  \qquad w \in \hilbH ,
\end{equation}
so that \(\ketbra{u}{v}^\adjoint = \ketbra{v}{u}\) and
\(T\ketbra{u}{v}S = \ketbra{Tu}{S^\adjoint v}\) for bounded \(T, S\).
Sesquilinear forms built from a map \(\varphi\), such as \(\ipphi{\cdot}{\cdot}\),
follow the comma convention.

% ---------------------------------------------------------------------------
\subsubsection{Linear maps}
\label{sec:conv-maps}

Let \(\varphi : \alg \to \bdd{\hilbH}\) be linear. We write
\(\varphi_n \eqdef \id_{\Mn} \tensor \varphi :
\Mn(\alg) \to \Mn(\bdd{\hilbH})\), and call \(\varphi\)
\emph{completely positive} if every \(\varphi_n\) is positive. We set
\begin{equation}
\label{eq:sharp}
  \varphi^\sharp(a) \eqdef \varphi(a^\adjoint)^\adjoint , \qquad a \in \alg .
\end{equation}
The map \(\varphi \mapsto \varphi^\sharp\) is a conjugate-linear involution on
the linear maps \(\alg \to \bdd{\hilbH}\). We call \(\varphi\) \emph{Hermitian}
if \(\varphi^\sharp = \varphi\), i.e.\ if \(\varphi\) sends self-adjoint elements
to self-adjoint operators. Every \(\varphi\) splits uniquely as
\(\varphi = \varphi_1 + i\varphi_2\) with
\(\varphi_1 = \frac12(\varphi + \varphi^\sharp)\) and
\(\varphi_2 = \frac1{2i}(\varphi - \varphi^\sharp)\) Hermitian.

\begin{remark}
The symbol \(\sharp\) is already in use elsewhere, and we keep it deliberately.
\begin{enumerate}
  \item \emph{Modular theory.} In Tomita--Takesaki theory \(\xi \mapsto \xi^\sharp\)
    is the involution of a left Hilbert algebra; for a von Neumann algebra
    \(M\) with cyclic separating vector \(\Omega\) it is
    \((x\Omega)^\sharp = x^\adjoint\Omega\), whose closure is the Tomita operator
    \(S\). Our use extends this one. Take \(\hilbH = \complex\) and a faithful
    tracial state \(\tau\) on \(\alg\), and write a functional as
    \(f = \tau(h\,\cdot\,)\). Then
    \[f^\sharp(a) = \overline{\tau(ha^\adjoint)} 
    = \tau(ah^\adjoint) = \tau(h^\adjoint a),\]
    so \(f^\sharp = \tau(h^\adjoint\,\cdot\,)\).
    On densities, \(\sharp\) is exactly the Tomita involution
    \(\vect h \mapsto \vect{h^\adjoint}\) of \(L^2(\alg, \tau)\) (here \(S = J\),
    since \(\tau\) is a trace). In both settings \(\sharp\) is the
    conjugate-linear involution whose fixed points are the Hermitian objects,
    and under \(\varphi \mapsto \Fphi\) it becomes the operator adjoint.
    \todo{cref the theorem \(F_{\varphi^\sharp} = \Fphi^\adjoint\) once ported.}
  \item \emph{Alternatives.} Some authors write \(\varphi^\adjoint\) for
    \eqref{eq:sharp}. We avoid it because \(\varphi^\adjoint\) also denotes the
    Hilbert--Schmidt (Banach-space) adjoint, which we use to pass between
    unital and trace-preserving maps.
    \todo{Confirm which references (e.g.\ Paulsen) use \(\varphi^\adjoint\).}
  \item \emph{Krein spaces.} \(T^\sharp\) often denotes the adjoint of an
    operator on a Krein space. We never write it that way; the Krein-space
    adjoint is named in words where it appears.
\end{enumerate}
\end{remark}

% ---------------------------------------------------------------------------
\subsection{The GNS representation}
\label{sec:gns}

\begin{theorem}[Gelfand--Naimark--Segal~{\cite[\S17, Theorem~2]{Naimark1972}}]
\label{thm:gns}
Let \(\alg\) be a unital C\(^*\)-algebra and \(\omega\) a positive linear
functional on \(\alg\). There are a Hilbert space \(\hilbH_\omega\), a
\(*\)-representation \(\pi_\omega\) of \(\alg\) on \(\hilbH_\omega\), and a vector
\(\Omega_\omega \in \hilbH_\omega\), cyclic for \(\pi_\omega(\alg)\), with
\begin{equation}
\label{eq:gns}
  \omega(a) = \ip{\pi_\omega(a)\Omega_\omega}{\Omega_\omega}
  = \bra{\Omega_\omega}\pi_\omega(a)\ket{\Omega_\omega} ,
  \qquad a \in \alg .
\end{equation}
The triple \((\pi_\omega, \Omega_\omega, \hilbH_\omega)\) is unique up to a
unitary intertwining the representations and the cyclic vectors.
\end{theorem}

Naimark writes \(\xi_0\) for the cyclic vector, though we write \(\Omega_\omega\)
as is common in the operator-algebra and physics literature.

% ---------------------------------------------------------------------------
\subsection{Stinespring's theorem}
\label{sec:stinespring}

Stinespring's theorem is the GNS construction for completely positive maps
into \(\bdd{\hilbH}\), with \(\hilbH\) in place of \(\complex\).

\begin{theorem}[Stinespring~{\cite[Theorem~1]{Stinespring1955}}]\label{thm:stinespring}
Let \(\alg\) be a unital C\(^*\)-algebra and \(\varphi : \alg \to \bdd{\hilbH}\)
completely positive. There are a Hilbert space \(\hilbK\), a \(*\)-representation
\(\pi\) of \(\alg\) on \(\hilbK\), and a bounded \(V : \hilbH \to \hilbK\) with
\begin{equation}
\label{eq:stinespring}
  \varphi(a) = V^\adjoint\pi(a)V , \qquad a \in \alg .
\end{equation}
\end{theorem}

The dilation can be taken with \(\pi\) unital, as in the construction below.
By~\cite[Theorem~4.1]{Paulsen2002}, we have additionally
\(V^\adjoint V = \varphi(1)\), so
\(\norm{V}^2 = \norm{V^\adjoint V} = \norm{\varphi(1)}\).
\todo{Check the theorem number in Paulsen.}

Such a triple \((\pi, V, \hilbK)\) is a \emph{Stinespring dilation} of
\(\varphi\). It is \emph{minimal} if \(\pi(\alg)V\hilbH\) spans a dense subspace
of \(\hilbK\). Minimal dilations are unique up to a unitary intertwining the
representations and the maps \(V\)~\cite{Arveson1969},
\cite[Chapter~4]{Paulsen2002}; this is not in Stinespring's paper.
\todo{Locate the precise statement in Paulsen, Chapter~4.}

We give a brief overview of the construction, as we will modify the presentation
slightly to obtain later results.

\paragraph{The construction.} Stinespring's proof follows GNS with
\(\hilbH \algtensor \alg\) in place of \(\alg\). Put
\begin{equation}
\label{eq:stine-sesqui-form}
  \ipphi{x \tensor a}{y \tensor b} \eqdef \ip{\varphi(b^\adjoint a)x}{y}
\end{equation}
on \(\hilbH \algtensor \alg\), extended sesquilinearly. For
\(u = \sum_{i=1}^m x_i \tensor a_i\),
\[
  \ipphi{u}{u} = \sum_{i,j} \ip{\varphi(a_j^\adjoint a_i)x_i}{x_j}
  = \ip{\varphi_m\br{[a_j^\adjoint a_i]_{j,i}}\vec x}{\vec x} \ge 0 ,
\]
since \(\varphi_m\) is positive
and the matrix \([a_j^\adjoint a_i]_{j,i}\) is positive in \(\Mn[m](\alg)\):
\[
  [a_j^\adjoint a_i]_{j,i} = X^\adjoint X, \qquad
  X \eqdef [\delta_{j,1} a_i]_{j,i} =
  \begin{pmatrix} a_1 & \cdots & a_m \\ 0 & \cdots & 0 \\ \vdots & & \vdots \end{pmatrix}.
\]
So~\eqref{eq:stine-sesqui-form} is positive semidefinite.
Note that is the one place complete positivity is used.
To make this a genuine inner product, let
\[
  \hilbN \eqdef \set{v \in \hilbH \algtensor \alg : \ipphi{v}{v} = 0} .
\]
Note that \(\hilbN\) is a
subspace, and \(\ipphi{\cdot}{\cdot}\) descends to an inner product on the
quotient \((\hilbH \algtensor \alg)/\hilbN\).
Let \(\hilbK\) be its completion, write \([v]\) for the class of \(v\), and set
\(\pi(c)[x \tensor a] \eqdef [x \tensor ca]\)
and \(Vx \eqdef [x \tensor 1]\). Then
\(\ip{V^\adjoint\pi(a)Vx}{y} = \ipphi{x \tensor a}{y \tensor 1}
= \ip{\varphi(a)x}{y}\), which is~\eqref{eq:stinespring}, and the dilation is
minimal by construction.

\begin{remark}
For \(\hilbH = \complex\) we have \(\complex \algtensor \alg = \alg\),
\eqref{eq:stine-sesqui-form} becomes \(\ip{a}{b}_\omega = \omega(b^\adjoint a)\),
\(V \cong \Omega_\omega\),%
\footnote{A bounded map \(V : \complex \to \hilbK\) is determined by \(V1\), and
\(V \mapsto V1\) identifies \(\bdd{\complex, \hilbK}\) with \(\hilbK\). So we may
write \(V = \ket{\Omega_\omega}\) and \(V^\adjoint = \bra{\Omega_\omega}\), and
then~\eqref{eq:stinespring} is~\eqref{eq:gns}:
\(V^\adjoint\pi(a)V = \bra{\Omega_\omega}\pi(a)\ket{\Omega_\omega}\).}
and the construction is the GNS representation of
\(\omega = \varphi\) (\cref{thm:gns}). Uniqueness of
minimal dilations specializes to uniqueness of the GNS representation, and
minimality to cyclicity of \(\Omega_\omega\). 
% More generally, for a unit
% vector \(x \in \hilbH\) the cyclic subspace \(\overline{\pi(\alg)Vx}\) carries
% the GNS representation of \(\ip{\varphi(\cdot)x}{x}\), and a dilation is
% minimal exactly when these subspaces span a dense subspace of \(\hilbK\).
\end{remark}

Working in finite dimensions allows us to replace the quotient-and-completion step.
%  is what this paper replaces.
Later, \todo{cref the section defining \(\Fphi\)} we fix the space
\(\hilbH \tensor L^2(\alg, \tau)\) in advance and realize~\eqref{eq:stine-sesqui-form} there
by a single operator.

% ---------------------------------------------------------------------------
\subsection{Commutants and opposite algebras}
\label{sec:commutants}

For a set \(S\) of bounded operators, \(S'\) denotes its commutant. The
\emph{opposite algebra} \(\alg\op\) is the vector space \(\alg\) with the same
involution and the reversed product
\begin{equation}
\label{eq:op}
  a\op b\op \eqdef (ba)\op ,
\end{equation}
where \(a\op\) denotes \(a\) regarded as an element of \(\alg\op\). Thus
\(a \mapsto a\op\) is a linear, \(*\)-preserving anti-isomorphism
\(\alg \to \alg\op\).

% Transpose remark: keep only if used later; otherwise delete.
\begin{remark}
For \(\alg = \Mn\) with its standard matrix units, the transpose gives a
\emph{linear} \(*\)-isomorphism
\begin{equation}
\label{eq:transpose-op}
  \Mn \to \Mn\op , \qquad a \mapsto (a^\transpose)\op .
\end{equation}
It is linear and bijective, multiplicative since
\(((ab)^\transpose)\op = (b^\transpose a^\transpose)\op
= (a^\transpose)\op (b^\transpose)\op\) by~\eqref{eq:op}, and \(*\)-preserving
since \((a^\adjoint)^\transpose = \overline{a} = (a^\transpose)^\adjoint\), with
\(\overline{a}\) the entrywise conjugate. 
% Applied blockwise it gives a
% \(*\)-isomorphism \(\alg \cong \alg\op\) for every finite-dimensional
% \(\alg \cong \bigoplus_k \Mn[n_k]\). Two caveats.
% \begin{enumerate}
%   \item It depends on the matrix units. Transposing in the basis
%     \(\set{ue_i}\), for a unitary \(u\), gives
%     \(u(u^\adjoint a u)^\transpose u^\adjoint
%     = \operatorname{Ad}(uu^\transpose)(a^\transpose)\). So the isomorphism is
%     canonical only up to an inner automorphism, and unchanged exactly when
%     \(uu^\transpose\) is scalar (e.g.\ \(u\) real orthogonal).
%   \item Beyond finite dimensions it has no analogue in general. The map
%     \(a \mapsto (a^\adjoint)\op\) is always a canonical \emph{conjugate-linear}
%     \(*\)-isomorphism \(\alg \to \alg\op\), but a linear one need not exist:
%     Connes gave a factor not anti-isomorphic to itself, and Phillips a simple
%     separable C\(^*\)-algebra not isomorphic to its opposite.
%     \todo{Add bib entries: Connes, Ann.\ Math.\ 101 (1975); Phillips, Proc.\
%     AMS 132 (2004).}
% \end{enumerate}
\end{remark}

% ---------------------------------------------------------------------------
\subsection{The standard form of a finite-dimensional algebra}
\label{sec:standard-form}

From now on \(\alg\) is a finite-dimensional C\(^*\)-algebra. By the structure
theorem~\cite[Theorem~III.1.1]{Davidson1996}, \cite{Takesaki1979},
\begin{equation}
\label{eq:structure}
  \alg \cong \bigoplus_{k=1}^r \Mn[n_k], \qquad a = (a_1, \dots, a_r).
\end{equation}
\todo{Check the theorem number in Davidson and add a precise reference in
Takesaki, Vol.~I.}
Let \(\iota_k\) and \(p_k\) be the inclusion of, and projection onto, the
\(k\)-th summand, and set \(\varphi_k \eqdef \varphi \compose \iota_k\). Then
\(\varphi = \sum_k \varphi_k \compose p_k\). Since \(\iota_k\) and \(p_k\) are
\(*\)-homomorphisms, \(\varphi\) is completely positive if and only if every
\(\varphi_k\) is. Dilations split the same way: if
\((\pi_k, V_k, \hilbK_k)\) is a minimal Stinespring dilation of \(\varphi_k\),
then \(\pi \eqdef \bigoplus_k \pi_k \compose p_k\) on
\(\hilbK \eqdef \bigoplus_k \hilbK_k\), with \(Vx \eqdef (V_kx)_k\), is a
minimal Stinespring dilation of \(\varphi\). It therefore suffices, for the
treatment of completely positive maps, to consider \(\alg = \Mn\).

% ---------------------------------------------------------------------------
\subsection{The space \texorpdfstring{\(\Ltwo\)}{L2(A)}}
\label{sec:l2-space}

Let \(\alg\) be a finite-dimensional C\(^*\)-algebra, as in
\cref{sec:standard-form}, with a faithful tracial state \(\tau\).
Following Anantharaman and Popa~\cite{AnantharamanPopa}, we take
\(\Ltwo = L^2(\alg, \tau)\) to be the GNS representation of \(\tau\)
(\cref{thm:gns}). \todo{Add the section number in Anantharaman--Popa where
\(L^2(M, \tau)\) is defined.} Since \(\tau\) is faithful and \(\dim \alg < \infty\), no
quotient or completion is needed: \(\Ltwo\) is \(\alg\) itself, written
\(a \mapsto \vect a\), with inner product
\begin{equation}
\label{eq:l2-ip}
  \ip{\vect a}{\vect b} \eqdef \tau(b^\adjoint a) ,
\end{equation}
linear in the first argument as in \cref{sec:conv-braket}, and cyclic vector
\(\Omega \eqdef \vect 1 = \Omega_\tau\). When we take \(\alg = \Mn\) to be a
matrix algebra and \(\tau\) to be the trace, the norm
\(\norm{\vect a}_2 = \tau(a^\adjoint a)^{1/2}\) is the Hilbert--Schmidt norm,
normalized by \(n^{-1/2}\) since \(\tau = \frac1n\trace\).

\paragraph{The standard representations.}
For \(a, b \in \alg\), let \(L_a\) and \(R_b\) be left and right multiplication
on \(\Ltwo\):
\begin{equation}
\label{eq:l2-mult}
  L_a\vect c \eqdef \vect{ac}, \qquad R_b\vect c \eqdef \vect{cb},
  \qquad c \in \alg .
\end{equation}
The \emph{standard representation} is
\[
  \lambda : \alg \to \bdd{\Ltwo}, \qquad \lambda(a) \eqdef L_a ,
\]
and it is the GNS representation of \(\tau\): \(\lambda = \pi_\tau\) and
\(\tau(a) = \ip{\lambda(a)\Omega}{\Omega}\). Right multiplication reverses
products, \(R_bR_c = R_{cb}\), so \(a \mapsto R_a\) is a representation of the
opposite algebra. The \emph{right standard representation} is
\[
  \rho : \alg\op \to \bdd{\Ltwo}, \qquad \rho(a\op) \eqdef R_a ,
\]
which is multiplicative by~\eqref{eq:op}:
\(\rho(b\op)\rho(c\op) = R_{cb} = \rho((cb)\op) = \rho(b\op c\op)\). 
Additionally, the two commute as
\(\lambda(a)\rho(b\op) = \rho(b\op)\lambda(a)\) by associativity.
\(\Ltwo\) is also known as the 
\emph{trivial} \(\alg\)--\(\alg\)Hilbert bimodule~\cite{AnantharamanPopa},
as \(a \cdot \xi \cdot b \eqdef \lambda(a)\rho(b\op)\xi\).

\begin{lemma}
\label{lemma:standard-rep}\leavevmode
\begin{enumerate}[label=(\roman*)]
  \item \(\lambda\) is a faithful unital \(*\)-representation of \(\alg\), and
    \(\lambda(\alg)\Omega = \Ltwo\).
  \item \(\rho\) is a faithful unital \(*\)-representation of \(\alg\op\), and
    \(\rho(\alg\op)\Omega = \Ltwo\).
  \item \(\lambda(\alg)' = \rho(\alg\op)\).
  \item For any Hilbert space \(\hilbH\),
    \((\id_{\hilbH} \tensor \lambda(\alg))' = \bdd{\hilbH} \tensor \rho(\alg\op)\),
    an algebraic tensor product that is already a von Neumann algebra.
\end{enumerate}
\end{lemma}
\begin{proof}
(i)--(iii) are standard; see~\cite{AnantharamanPopa}. \todo{Add the section
number in Anantharaman--Popa.}
(iv) This follows from the commutation theorem for tensor products of von
Neumann algebras, \((M \mathbin{\bar\tensor} N)' = M' \mathbin{\bar\tensor} N'\),
where \(\bar\tensor\) is the weak closure of the algebraic tensor
product~\cite[Theorem~IV.5.9]{Takesaki1979}. \todo{Check the theorem number in
Takesaki.} Take \(M = \complex\id_{\hilbH}\) and \(N = \lambda(\alg)\), so that
\(M' = \bdd{\hilbH}\) and \(N' = \rho(\alg\op)\) by~(iii). Since \(\alg\) is
finite-dimensional, \(\bdd{\hilbH} \tensor \rho(\alg\op)\) is already weakly
closed so no closure is needed.
\end{proof}

\bibliographystyle{plain}
\bibliography{references}

% Draft only: remove before submission.
\listoftodos

\end{document}
